
\chapter{Methodology}
\label{chap:method}

In this chapter we describe the three-component framework used to reduce FEM
solver calls in a transient water-quenching simulation.
The \textbf{FEM Reference Solver} (Section~\ref{sec:mswp}) computes the
ground-truth temperature field at selected anchor time steps using a
Crank--Nicolson finite-difference scheme; it is called only at anchor points
and the NAS-PINN replaces all intermediate steps.
The \textbf{NAS-PINN} (Sections~\ref{sec:pinn_form}--\ref{sec:nas_method}) is
a feed-forward neural network trained independently for each prediction window;
its architecture is selected automatically by Neural Architecture Search before
training begins, and its output is constrained to match the FEM anchor field
exactly at the window start.
The \textbf{MSWP Controller} (Section~\ref{sec:mswp}) decides when to call FEM
and when to call the NAS-PINN via the skip factor $k$: at $k=1$ FEM anchors
every step, at $k=5$ the NAS-PINN covers five steps per FEM call saving 80\%
of solver invocations, and an adaptive-$k$ variant adjusts this factor at run
time.
Section~\ref{sec:mortensen_reference} defines the target physical problem,
Section~\ref{sec:domains} introduces the benchmark domains, and
Section~\ref{sec:metrics} defines the evaluation metrics.

% =============================================================================
\section{Reference Problem: Mortensen et al.\ Industrial Simulation (A356 Quenching)}
\label{sec:mortensen_reference}

The physical problem this thesis targets is the transient water quenching
of A356 aluminium automotive subframes studied by Mortensen et al.\
\cite{mortensen2026}.
This section describes their simulation, the physical parameters we adopted,
and the simplifications we made.

\subsection{Industrial Context and Objective}

Mortensen et al.\ address a core challenge in aluminium automotive
manufacturing.
Cast subframes undergo non-uniform cooling during water quenching, generating
temperature gradients that drive residual stress and geometric distortion.
Dimensional tolerances are tight.
Deviations at key assembly points must stay within a few hundred microns.
Their work predicts these distortions using FEM and shows that rack
configuration changes can reduce them significantly.

\subsection{The Mortensen et al. FEM Model}

Mortensen et al.\ use StaMiSim, an in-house thermo-elastoviscoplastic FEM
solver developed at IFE (Institute for Energy Technology).
It simulates the full process chain from casting to final quenching.
The simulation covers three sequential phases.
First, as-cast (AC) deformation models solidification inside the mould,
including mould thermal expansion, airgap formation, and degating.
The AC simulation uses approximately 3~million finite elements across 27 coupled
domains representing the part and surrounding equipment.
Second, heat treatment creep holds the part at $540\,^\circ\mathrm{C}$ for
several hours.
Creep under gravity causes additional deformation depending on rack positioning.
Third, water quenching submerges the part in a water tank at approximately
$20\,^\circ\mathrm{C}$.
This is the phase this thesis targets.

The quenching thermal model uses a full temperature-dependent boiling curve
$h(T_s)$ comprising five heat transfer regimes, ranging from film boiling at
high surface temperatures to single-phase convection at low temperatures.
The effective heat transfer coefficient varies by orders of magnitude over the
30-second quench.
The model additionally accounts for vapour accumulation, orientation-dependent
cooling, and cavity detection via ray-beam analysis on the FE mesh.
Validated against coordinate measuring machine (CMM) measurements from
industrial casthouse trials, the model shows that rack optimisation reduces
peak distortions from ${\approx}3\,\mathrm{mm}$ to ${\approx}1\,\mathrm{mm}$.
The steepest temperature gradients occur in the first two seconds of quenching
and are the primary driver of bending distortion.

\subsection{What This Thesis Takes from Mortensen et al.}

This framework targets the thermal sub-problem of the Mortensen simulation,
that is, the transient temperature field during water quenching.
Accurate temperature prediction is the necessary prerequisite for the
subsequent mechanical distortion calculation.

\begin{table}[H]
\centering
\caption{Comparison of the Mortensen et al.~\cite{mortensen2026} industrial
framework and the NAS-PINN framework developed in this thesis.}
\label{tab:mortensen_comparison}
\small
\begin{tabular}{p{3.8cm}p{5.0cm}p{5.0cm}}
\toprule
\textbf{Element} & \textbf{Mortensen et al.~\cite{mortensen2026}} & \textbf{This thesis} \\
\midrule
\multicolumn{3}{l}{\textit{\textbf{Taken directly from Mortensen et al.}}} \\
\midrule
Material & A356 aluminium alloy & A356 aluminium alloy \\
Initial temperature & $T_0 = 540\,^\circ\mathrm{C}$ & $T_0 = 540\,^\circ\mathrm{C}$ \\
Quench temperature & $T_w = 20\,^\circ\mathrm{C}$ & $T_w = 20\,^\circ\mathrm{C}$ \\
Quench duration & 30 seconds & 30 seconds \\
Thermal diffusivity & $\alpha = k_T/(\rho c_p)$ & Same value \\
Engineering tolerance & Distortion $\lesssim$ few hundred $\mu$m & Temperature MAE $\leq 10$--$15\,^\circ\mathrm{C}$ \\
Process motivation & Rack optimisation reduces distortion & Accelerate the thermal sub-problem \\
\midrule
\multicolumn{3}{l}{\textit{\textbf{Simplifications relative to Mortensen et al.}}} \\
\midrule
Physical problem & Thermo-elastoviscoplastic (thermal + mechanical coupled) & Thermal only (heat equation) \\
Solver & StaMiSim (in-house FEM, multi-physics) & Custom Crank--Nicolson FD + NAS-PINN \\
Mesh size & $\approx 3\times10^6$ finite elements & $\leq 1\,300$ structured grid nodes \\
Domain count & 27--30 coupled domains (mould, rack, part\ldots) & 1 domain per benchmark \\
Time-step strategy & Adaptive $\Delta t$, implicit solver & Fixed $\Delta t_\mathrm{int}=0.5\,\mathrm{s}$ \\
Heat transfer $h$ & Full 5-regime boiling curve $h(T)$; spatial variation by rack position and cavity detection & Constant $h = 4\,000\,\mathrm{W/(m^2K)}$ \\
Material properties & Temperature-dependent $k_T(T)$, $\rho c_p(T)$; viscoplastic flow stress & Constant $k_T = 160\,\mathrm{W/(mK)}$, $\rho c_p = 2.4\times10^6\,\mathrm{J/(m^3K)}$ \\
Geometry & Proprietary CAD: A356 subframe with rails, enclosed cavities, support pads & Canonical 2D/3D geometries + 3D Thermal Fin benchmark \\
Cavity treatment & Ray-beam vapour detection; position-dependent correction & Uniform BC; cavity modelled as low-$h$ region \\
Validation & Physical CMM measurements on production parts & Comparison to own FEM reference solution \\
\bottomrule
\end{tabular}
\end{table}

Table~\ref{tab:mortensen_comparison} shows precisely what is taken and what
is simplified.
The physical parameters and engineering tolerances that define the validation
target of this thesis are taken directly from~\cite{mortensen2026}: A356
aluminium, $T_0=540\,^\circ\mathrm{C}$, $T_w=20\,^\circ\mathrm{C}$,
temperature accuracy within $10$--$15\,^\circ\mathrm{C}$.
The five-regime boiling curve, temperature-dependent material properties, and
the real industrial mesh are simplified to allow the MSWP framework to be
validated systematically before being applied to the full model.
The NAS-PINN framework is designed so that its thermal field predictions can
serve as input to the mechanical distortion model of~\cite{mortensen2026} in
the complete pipeline.
The extension path is described in Section~\ref{sec:conc_future}.

% =============================================================================
\section{Multi-Step Window Prediction (MSWP) Framework}
\label{sec:mswp}

The MSWP concept and its three operating modes are introduced in
Chapter~\ref{chap:intro}.
This section defines the prediction loop formally using three algorithms and
specifies the four experimental configurations.

\textbf{How the three components interact.}
At the start of window $i$, the FEM solver provides the current temperature
field $T^\mathrm{FEM}_{t_i}$ as the anchor input.
The NAS-PINN uses this field as its initial condition and predicts the
temperature over the next $k$ intervals (see Section~\ref{sec:pinn_form} for
the formulation; Section~\ref{sec:nas_method} for the selected architectures).
The MSWP controller records the per-window prediction error and uses it to
decide whether to keep, increase, or decrease $k$ for the next window.
This three-way handshake repeats until $t = 30\,\mathrm{s}$.

Figure~\ref{fig:naspinn_process} shows the full process from problem setup
through NAS architecture selection to evaluation.

\begin{figure}[htbp]
  \centering
  \includegraphics[width=\textwidth]{fig6_naspinn_process.png}
  \caption{NAS-PINN process flow from problem setup to evaluation.}
  \label{fig:naspinn_process}
\end{figure}

\subsection{FEM Reference Solver (Crank--Nicolson, Structured Grid)}
\label{subsec:fem_solver}

We used a Crank--Nicolson finite-difference solver on structured grids as
the FEM reference throughout this thesis.
The solver implements the transient heat equation:

\begin{equation}
\rho c_p \frac{\partial T}{\partial t} = \nabla \cdot (k_T \nabla T),
\label{eq:heat_dim}
\end{equation}

with a Robin boundary condition $k_T \partial T / \partial n = -h(T - T_w)$
on all exposed surfaces and constant material properties
($k_T = 160\,\mathrm{W/(mK)}$, $\rho c_p = 2.4 \times 10^6\,\mathrm{J/(m^3K)}$).

The Crank--Nicolson scheme is unconditionally stable and uses an internal
time step of $\Delta t_\mathrm{int} = 0.5\,\mathrm{s}$.
We exported (recorded) the full temperature field every $\Delta t_a = 1.5\,\mathrm{s}$,
producing 21 snapshots over the 30-second quench.
These exported snapshots serve as the \emph{anchor fields}: the FEM ground
truth that the NAS-PINN is trained to match.
The internal integration continues at $0.5\,\mathrm{s}$ steps between
exports; the PINN never modifies the FEM state.

Figure~\ref{fig:fem_solver_method} illustrates the full FEM workflow for
a representative 2D domain: domain discretisation, element assembly, global
system solution, and the comparison of the FEM reference field against the
NAS-PINN prediction.

\begin{figure}[H]
\centering
\includegraphics[width=\textwidth]{FEM_Proses.png}
\caption{FEM reference solver workflow (Circle 2D benchmark).
Panel~1: transient temperature field at $t=3\,\mathrm{s}$ obtained from the
FEM solver.
Panel~2: domain discretisation into triangular finite elements (nodes shown
as blue dots; elements as triangles).
Panel~3: element-level matrix assembly using material properties and shape
functions.
Panel~4: global conductivity $\mathbf{K}$, capacitance $\mathbf{C}$, and load
$\mathbf{f}$ matrices assembled; time step solved by
$(C + \Delta t K)T^{n+1} = \Delta t F^{n+1} + C T^n$.
Panel~5: absolute prediction error between the NAS-PINN and the FEM reference
at $t=3\,\mathrm{s}$, $k=1$ ($L_2 = 0.0294$).}
\label{fig:fem_solver_method}
\end{figure}

\subsection{Skip Factor and Anchor Reduction (FEM Call Budget)}

FEM temperature fields are exported every $\Delta t_a = 1.5\,\mathrm{s}$,
giving 20 anchor intervals across the full 30-second simulation.
The skip factor $k$ controls how many of those intervals one PINN window
covers at a time:

\begin{equation}
\Delta t_k = k\,\Delta t_a.
\label{eq:dtk}
\end{equation}

At $k=1$ the network predicts one interval and FEM is still called at every
step, so no FEM calls are eliminated.
As $k$ grows, each PINN window gets longer and fewer FEM calls are needed.
The total number of windows and the corresponding anchor reduction are:

\begin{equation}
N_\mathrm{win}(k) = \left\lceil \frac{20}{k} \right\rceil, \qquad
R(k) = \left(1 - \frac{N_\mathrm{win}(k)}{20}\right)\times 100\%.
\label{eq:reduction}
\end{equation}

Table~\ref{tab:k_structure} shows what this looks like in practice for
each value of $k$ tested.

\begin{table}[H]
\centering
\caption{FEM anchor structure for each skip factor $k$.}
\label{tab:k_structure}
\begin{tabular}{ccccc}
\toprule
$k$ & Anchor intervals & Window length & No.\ of windows & Anchor reduction \\
\midrule
1 & 1 & $1.5\,\mathrm{s}$ & 20 & $0\%$ \\
2 & 2 & $3.0\,\mathrm{s}$ & 10 & $50\%$ \\
3 & 3 & $4.5\,\mathrm{s}$ &  7 & $65\%$ \\
4 & 4 & $6.0\,\mathrm{s}$ &  5 & $75\%$ \\
5 & 5 & $7.5\,\mathrm{s}$ &  4 & $80\%$ \\
\bottomrule
\end{tabular}
\end{table}

The three operating modes for different skip factors are illustrated in
Figure~\ref{fig:timestep_intro} in Chapter~\ref{chap:intro}.

\subsection{Operating Modes (Fixed-$k$, Adaptive-$k$, PINN-only)}
\label{subsec:adaptive_k}

Experiments were run in two $k$ selection modes.
In fixed $k$ mode, the skip factor is set once before the simulation and stays
constant for every window.
We swept $k=1,2,3,4,5$ for all three architectures on all domains, giving
15 runs per domain group.
This mode applies to both the FEM-anchored hybrid and the PINN-only rollout.

Algorithms~\ref{alg:fem_fixed} and~\ref{alg:pinn_only} show the two fixed $k$
rollout protocols.
The FEM-anchored version tests how much FEM work can be removed while still
correcting accumulated error at every anchor.
The PINN-only version is a stress test: it shows the largest apparent FEM
saving, but also exposes how quickly errors grow when the FEM correction is
removed.

\begin{algorithm}[H]
\caption{FEM-anchored fixed $k$ MSWP}
\label{alg:fem_fixed}
\begin{algorithmic}[1]
\Require fixed skip factor $k$,\; stored FEM snapshots
\For{each prediction window $i$}
    \State Set $\theta_\mathrm{ic} \leftarrow T^\mathrm{FEM}_{t_i}$
    \State Train the IC consistent PINN over $[t_i,t_i+k\Delta t_a]$
    \State Predict $\hat{T}_{t_{i+1}}$ at the window endpoint
    \State Compute the window error against $T^\mathrm{FEM}_{t_{i+1}}$
    \State Use the next FEM snapshot as the IC for the following window
\EndFor
\State Report the mean window error and FEM anchor reduction
\end{algorithmic}
\end{algorithm}

\begin{algorithm}[H]
\caption{PINN-only fixed $k$ rollout}
\label{alg:pinn_only}
\begin{algorithmic}[1]
\Require fixed skip factor $k$,\; initial FEM snapshot $T^\mathrm{FEM}_{0}$
\State Set $\theta_\mathrm{ic} \leftarrow T^\mathrm{FEM}_{0}$
\For{each prediction window $i$}
    \State Train the IC consistent PINN over $[t_i,t_i+k\Delta t_a]$
    \State Predict $\hat{T}_{t_{i+1}}$ at the window endpoint
    \State Compute the window error against the stored FEM reference
    \State Set $\theta_\mathrm{ic} \leftarrow \hat{T}_{t_{i+1}}$
    \Comment{no FEM correction}
\EndFor
\State Report accumulated rollout error and the 95\% FEM call reduction
\end{algorithmic}
\end{algorithm}

In adaptive $k$ mode, the controller adjusts $k$ after each window based on
how accurately the PINN predicted that window.
If the error stayed below a promotion threshold for two consecutive windows,
$k$ is increased by one.
If it exceeded a demotion threshold, $k$ drops back by one.
Phase-based caps prevent the controller from using long windows during the
high-gradient early transient.
This mode was applied in two configurations: with FEM anchoring and without
anchoring.
In the PINN-only adaptive variant, the IC for each window comes from the
previous PINN output and FEM is used only at $t=0$.

The window error $\varepsilon_i$ used by the controller is always computed by
comparing the PINN prediction at the end of window $i$ to the precomputed FEM
reference snapshot at the same time point.
In the FEM-anchored case this comparison is part of the regular anchor step.
In the PINN-only case the FEM reference is available because a full FEM
solution was run in advance.
The controller reads this stored snapshot to evaluate $\varepsilon_i$ without
triggering any additional FEM calls at deployment time.

Algorithm~\ref{alg:adaptive_k} gives the complete decision procedure used
in both adaptive configurations.

\begin{algorithm}[H]
\caption{Adaptive $k$ controller (used for FEM-anchored adaptive and
PINN-only adaptive modes)}
\label{alg:adaptive_k}
\begin{algorithmic}[1]
\Require $k \leftarrow k_\mathrm{init}$,\;
         thresholds $\tau_\mathrm{up}$, $\tau_\mathrm{down}$,\;
         streak required $s_\mathrm{req}$,\;
         phase caps $k_\mathrm{max}^\mathrm{early}$,\;
         $k_\mathrm{max}^\mathrm{mid}$,\;
         $k_\mathrm{max}^\mathrm{late}$,\;
         bounds $[k_\mathrm{min}, k_\mathrm{max}]$
\State $\mathit{streak} \leftarrow 0$
\For{each prediction window $i$}
    \State Run PINN over window $i$ with current $k$
    \State $\varepsilon_i \leftarrow \mathrm{MAE}(\hat{T}_i,\; T^\mathrm{FEM}_{t_i})$
    \Comment{compare to stored FEM snapshot}
    \If{$\varepsilon_i < \tau_\mathrm{up}$}
        \State $\mathit{streak} \leftarrow \mathit{streak} + 1$
        \If{$\mathit{streak} \geq s_\mathrm{req}$}
            \State $k \leftarrow \min(k+1,\;k_\mathrm{max}^\mathrm{phase},\;k_\mathrm{max})$
            \Comment{promote}
            \State $\mathit{streak} \leftarrow 0$
        \EndIf
    \ElsIf{$\varepsilon_i > \tau_\mathrm{down}$}
        \State $k \leftarrow \max(k-1,\;k_\mathrm{min})$
        \Comment{demote}
        \State $\mathit{streak} \leftarrow 0$
    \Else
        \State $\mathit{streak} \leftarrow 0$
        \Comment{hold}
    \EndIf
    \State \textbf{Set IC for next window:} FEM snapshot (FEM-anchored) or
    \Statex \hspace{4.8em} PINN output (PINN-only)
\EndFor
\end{algorithmic}
\end{algorithm}

The phase cap $k_\mathrm{max}^\mathrm{phase}$ is
$k_\mathrm{max}^\mathrm{early}=2$ for $t<10\,\mathrm{s}$,
$k_\mathrm{max}^\mathrm{mid}=4$ for $10\leq t<20\,\mathrm{s}$, and
$k_\mathrm{max}^\mathrm{late}=5$ after that.
For the Thermal Fin experiments we used $k_\mathrm{init}=1$,
$\tau_\mathrm{up}=13\,^\circ\mathrm{C}$,
$\tau_\mathrm{down}=16\,^\circ\mathrm{C}$, and $s_\mathrm{req}=2$.
Table~\ref{tab:modes} summarises all four configurations.
The three algorithms produce four operating modes because the adaptive-$k$
controller is shared between the FEM-anchored and PINN-only variants.

\begin{table}[H]
\centering
\caption{Experimental configurations: three algorithms produce four operating modes.}
\label{tab:modes}
\small\setlength{\tabcolsep}{4pt}
\begin{tabular}{@{}p{3.4cm}p{2.0cm}p{2.6cm}p{1.6cm}p{4.0cm}@{}}
\toprule
Mode & Algorithm & IC source & FEM saved & Domains \\
\midrule
FEM-anchored, fixed $k$
  & Alg.~\ref{alg:fem_fixed}
  & FEM snapshot
  & 0--80\%
  & All 8 canonical domains \\
FEM-anchored, adaptive $k$
  & Alg.~\ref{alg:adaptive_k}
  & FEM snapshot
  & $\sim$45--85\%
  & 2D adaptive, 3D canonical, Thermal Fin \\
PINN-only, fixed $k$
  & Alg.~\ref{alg:pinn_only}
  & Prev.\ PINN output
  & 95\%
  & Thermal Fin \\
PINN-only, adaptive $k$
  & Alg.~\ref{alg:adaptive_k}
  & Prev.\ PINN output
  & 95\%
  & Thermal Fin \\
\bottomrule
\end{tabular}
\end{table}

% ─────────────────────────────────────────────────────────────────────────────
\section{NAS-PINN Formulation (IC-Consistent Output, Loss, Stabilisation)}
\label{sec:pinn_form}

The theoretical background (spectral bias, Fourier features, self-adaptive
loss weighting) is in Chapter~\ref{chap:theory}.
This section covers the implementation decisions specific to the MSWP setting:
the IC-consistent output form, loss function, Fourier embedding, and collocation
strategy.
Figure~\ref{fig:pinn_framework} shows the complete single-window pipeline
from FEM anchor input through NAS architecture selection to the IC-consistent
output and physics-based loss terms.

\begin{figure}[htbp]
\centering
\includegraphics[width=\textwidth]{nas_pinn_framework.png}
\caption{Single-window NAS-PINN prediction pipeline.}
\label{fig:pinn_framework}
\end{figure}

The pipeline operates as follows.
At the start of each prediction window, the FEM solver provides the
temperature field $u_\mathrm{ic}$ as the anchor input.
The spatial coordinates $(x, y, z, \tau)$ and $u_\mathrm{ic}$ enter the
network together.
An optional Fourier feature transformation can replace the raw spatial
coordinates with a sinusoidal projection before the network layers; without
it, the coordinates are passed directly.
The NAS module selects one of three candidate architectures, Bayesian/TPE
($5{\times}151$, ReLU), NSGA-II ($3{\times}153$, tanh), or NSGA-III
($3{\times}75$, tanh), and fixes it for all windows of the simulation.
The network output is wrapped by an IC-consistent layer that enforces
$\hat{u}(\mathbf{x},0) = u_\mathrm{ic}$ exactly for every set of weights,
so the training loss never needs to fight the initial condition.
Four loss terms then supervise training: the PDE residual
$\mathcal{L}_\mathrm{pde}$, the Robin boundary condition
$\mathcal{L}_\mathrm{bc}$, the initial condition check
$\mathcal{L}_\mathrm{ic}$, and endpoint supervision
$\mathcal{L}_\mathrm{end}$ that penalises deviation from the next FEM
snapshot.
Their weighted sum is minimised by backpropagation.

\subsection{Network Input and Output (IC-Consistent Form)}

Within each window $[t_s, t_s+\Delta t_k]$, the local normalised time is
$\tau = (t-t_s)/\Delta t_k \in [0,1]$.
The network input is $(x, y, \tau, \theta_\mathrm{ic})$ in 2D and
$(x, y, z, \tau, \theta_\mathrm{ic})$ in 3D, where $\theta_\mathrm{ic}$ is
the dimensionless temperature at the window start.

The output uses the IC consistent form:

\begin{equation}
\hat\theta(\mathbf{x},\tau)
= \theta_\mathrm{ic}(\mathbf{x}) + \tau\,N_\phi(\mathbf{x},\tau,\theta_\mathrm{ic}),
\label{eq:ic_form}
\end{equation}

which guarantees $\hat\theta(\mathbf{x},0)=\theta_\mathrm{ic}(\mathbf{x})$
by construction \cite{lagaris1998}.
This removes the need for a separate IC residual loss term at $\tau=0$ and
allows the network to focus on learning the temperature change after the
anchor.

We evaluated Fourier feature embedding as a candidate input transformation.
Spatial coordinates are passed through a fixed random sinusoidal projection
before entering the MLP:
\begin{equation}
\gamma(\mathbf{x}) = \bigl[\cos(2\pi \mathbf{B}\mathbf{x}),\;
                           \sin(2\pi \mathbf{B}\mathbf{x})\bigr],
\qquad \mathbf{B} \sim \mathcal{N}(0,\sigma^2).
\label{eq:fourier_embed}
\end{equation}
A targeted ablation (Appendix~\ref{app:ablation}) showed that this embedding
reduces mean-window MAE by $1.8$--$4.4\,^\circ\mathrm{C}$ on the Thermal
Fin for Bayesian/TPE.
It is detrimental for the compact NSGA-II and NSGA-III tanh networks.
The failure of Fourier embedding on NSGA-II/III is attributable to
dimensionality mismatch: the embedding projects the $(d+2)$-dimensional input
through $2F=128$ sinusoidal features, quadrupling the effective input
dimension.
A three-layer network with 75 or 153 neurons per layer does not have
sufficient capacity to disentangle this high-dimensional representation while
simultaneously satisfying the PDE and boundary conditions.
Consequently, Fourier embedding is used only for Bayesian/TPE on the Thermal
Fin benchmark.
All other reported results use the plain IC-consistent MLP.

\textbf{Limitation on architecture comparison.}
The use of Fourier embedding for Bayesian/TPE but not for NSGA-II/III on the
Thermal Fin benchmark introduces a confound: the observed performance
differences partly reflect the embedding choice rather than the architecture
alone.
To isolate this effect, Appendix~\ref{app:ablation} reports a controlled
ablation comparing plain and Fourier-embedded Bayesian/TPE at matched epoch
budgets, confirming that the embedding accounts for $1.8$--$4.4\,^\circ\mathrm{C}$
of the gap.
A fully fair three-way architecture comparison on the Thermal Fin would require
an embedding scheme that benefits all three architectures equally.
This is deferred to future work.

\subsection{Loss Function (PDE, BC, IC, Endpoint Supervision)}

The total training loss for the FEM-anchored protocol is:

\begin{equation}
\mathcal{L}
= w_\mathrm{pde}\mathcal{L}_\mathrm{pde}
+ w_\mathrm{ic}\mathcal{L}_\mathrm{ic}
+ w_\mathrm{bc}\mathcal{L}_\mathrm{bc}
+ w_\mathrm{end}\mathcal{L}_\mathrm{end}.
\label{eq:loss_full}
\end{equation}

Each term is defined as:

\begin{align}
\mathcal{L}_\mathrm{pde}
&= \mathrm{MSE}\!\left(
   \frac{1}{\Delta t_k}\frac{\partial\hat\theta}{\partial\tau}
   - \alpha\nabla^2\hat\theta
   \right),
\label{eq:lpde2}\\[4pt]
\mathcal{L}_\mathrm{bc}
&= \mathrm{MSE}\!\left(
   \nabla\hat\theta\cdot\hat{\mathbf{n}}
   + \frac{h}{k_T}\hat\theta
   \right),
\label{eq:lbc2}\\[4pt]
\mathcal{L}_\mathrm{ic}
&= \mathrm{MSE}\!\left(
   \hat\theta(\mathbf{x},0) - \theta_\mathrm{ic}(\mathbf{x})
   \right),
\label{eq:lic2}\\[4pt]
\mathcal{L}_\mathrm{end}
&= \mathrm{MSE}\!\left(
   \hat\theta(\mathbf{x},1) - \theta_\mathrm{FEM}(\mathbf{x},t_s+\Delta t_k)
   \right).
\label{eq:lend2}
\end{align}

In the PINN-only protocol, $w_\mathrm{end}=0$ because the FEM field at the
window end is not available.

\subsection{Collocation Point Sampling (Interior and Boundary)}

Interior collocation points are sampled uniformly inside the domain by
rejection sampling.
For the Thermal Fin, a fraction of points is concentrated in challenging
subregions such as fin tips and reentrant corners.

Boundary collocation points are distributed proportionally to the surface area
of each face.
For the Thermal Fin benchmark, all 22 exposed faces receive area-proportional
samples.

\subsection{Stabilisation Choices}
\label{subsec:stabilisation_choices}

Several implementation choices were added because the plain windowed PINN
was not reliable enough on long windows and complex geometries.
Table~\ref{tab:stabilisation_choices} summarises what was used, why, and
where the effect appears in the results.
A quantitative improvement is stated only where a direct comparison was run;
otherwise the entry is reported as a design choice that supports the final
configuration.

\begin{table}[H]
\centering
\caption{Stabilisation choices used in the final MSWP-PINN framework, with
their literature basis and how they were integrated in this thesis.}
\label{tab:stabilisation_choices}
\scriptsize
\setlength{\tabcolsep}{3pt}
\begin{tabular}{p{0.16\textwidth}p{0.27\textwidth}p{0.27\textwidth}p{0.22\textwidth}}
\toprule
Choice & Literature basis & How we used it & Evidence in this thesis \\
\midrule
IC consistent output
& Lagaris et al.\ used trial solutions that satisfy constraints by
construction~\cite{lagaris1998}; PINNs then optimise the PDE and boundary
residuals around that constrained form~\cite{raissi2019}.
& We wrote each window prediction as
$\hat\theta=\theta_\mathrm{ic}+\tau N_\phi$, so the FEM anchor field is
matched exactly at $\tau=0$.
& Used in all reported runs.  The low $k=1$ errors in the 2D and 3D
benchmarks show that the anchored window formulation reproduces the FEM
trajectory before long window drift dominates. \\
\midrule
Fourier feature embedding
& Spectral bias makes standard networks learn low frequency components first
\cite{rahaman2019}; Fourier features improve the representation of
high frequency coordinate functions~\cite{tancik2020}.
& Evaluated as a candidate stabilisation choice.  The reported benchmark
results use the plain IC consistent MLP without Fourier embedding; a
targeted Fourier vs.\ no Fourier comparison at $k=2$ is reported in
Appendix~\ref{app:ablation}.
& See Appendix~\ref{app:ablation}.  Fourier features were available in the
implementation but were not found to be critical for the canonical
benchmark geometries within the tested training budget. \\
\midrule
Endpoint supervision and FEM anchoring
& Long time PINN rollouts are known to accumulate error and become difficult
to train~\cite{wang2021,krishnapriyan2021}; time decomposed PINN methods
reduce this by solving shorter subproblems~\cite{meng2020ppinn,wang2024causality}.
& We used the stored FEM field at the next anchor as an endpoint loss in
FEM-anchored mode.  In PINN-only mode this term is disabled, so the next
window starts only from the previous PINN output.
& Thermal Fin anchored runs stay near $8$--$16\,^\circ\mathrm{C}$, while
PINN-only rollouts rise to $59$--$217\,^\circ\mathrm{C}$. \\
\midrule
Adam followed by L-BFGS
& Adam provides robust stochastic optimisation~\cite{kingma2015}; L-BFGS is
a quasi Newton local optimiser often used for refinement in PINN
training~\cite{liu1989,byrd1995}.
& Each window was trained first with Adam and then refined with L-BFGS.  More
complex domains and compact NSGA networks used longer budgets.
& The extra epochs improve accuracy for NSGA-II/III but do not change the
optimal $k$ on most domains. \\
\midrule
Per-window re-initialisation
& Time decomposed PINN formulations reuse information between neighbouring
time blocks to improve long time training stability~\cite{meng2020ppinn,wang2024causality}.
& Each window begins from a new random network initialisation, ensuring
independence of successive training runs.
& Weight-transfer initialisation (carrying trained weights from the previous
window) was evaluated but provided no consistent accuracy benefit; per-window
re-initialisation was therefore adopted throughout. \\
\midrule
Feature balanced sampling
& PINNs rely on collocation points to enforce the PDE and boundary
residuals~\cite{raissi2019,karniadakis2021}; under sampled regions can
therefore dominate the error even if the global loss looks acceptable.
& We reserved fixed shares of collocation points for thin supports, cavities,
fin tips, and corner regions instead of sampling only by total volume.
& Used for the Thermal Fin case, where fin tips and corner regions
receive additional points. \\
\midrule
Seed sensitivity evaluation
& Stochastic neural network initialisation means results vary between runs;
quantifying this variance is necessary to assess whether reported single-run
values are representative.
& Each (domain, architecture, $k$) configuration was run with ten independent
random seeds.  All reported values are ten-seed means; standard deviations
and per-seed breakdowns are in the consolidated tables in each results section.
& Standard deviations remain below $1.4\,^\circ\mathrm{C}$ on canonical
benchmarks, confirming results are representative. \\
\bottomrule
\end{tabular}
\end{table}


\section{NAS-PINN Architecture Search (Bayesian/TPE, NSGA-II, NSGA-III)}
\label{sec:nas_method}

\subsection{What is NAS-PINN?}

Wang et al.~\cite{wang2023} introduced NAS-PINN as a framework in which
Neural Architecture Search selects the PINN architecture automatically.
Instead of manually tuning network depth, width, and activation, the NAS
module optimises these choices over a defined search space before training
begins on the target problem.

\textbf{Network structure.}
A NAS-PINN is a standard feed-forward MLP with a uniform hidden structure:
$L$ hidden layers, each with $W$ neurons and the same activation function $\sigma$.
All layers share the same width; only depth, width, and activation are searched.
The input is the tuple $(\mathbf{x}, \tau, \theta_\mathrm{ic})$: spatial
coordinates, normalised window time, and the FEM anchor temperature field.
The output is a scalar correction $\delta\theta$; the IC-consistent form wraps
it as $\hat\theta = \theta_\mathrm{ic} + \tau\,\delta\theta$ (Section~\ref{sec:pinn_form}).

We used the same search space as Wang et al.~\cite{wang2023}.
The number of hidden layers ranged over $L \in \{2,3,4,5,6\}$.
Neurons per layer ranged over $W \in \{32, 64, \ldots, 256\}$ with uniform
width across layers.
The activation function was selected from
$\sigma \in \{\text{ReLU},\, \text{tanh},\, \text{SiLU},\, \text{GELU}\}$.
The full search space table is in Appendix~\ref{app:nas_space}.

\textbf{NAS-once, deploy-everywhere.}
We ran each search once on the 2D rectangle domain at $k=1$.
The winning architecture was then fixed and used for all domains and skip
factors without repeating the search.
This strategy reflects industrial practice: re-running NAS for every geometry
is not affordable.
We treated architectural robustness under geometric transfer as a result to
be evaluated, not an assumption to be made.

\subsection{Search Strategies and Selected Architectures}

The theoretical background for each search strategy is in
Chapter~\ref{chap:theory}.
We compared three strategies representing different search philosophies:
single-objective accuracy maximisation (Bayesian/TPE), two-objective
accuracy-compactness trade-off (NSGA-II), and diversity-enforced Pareto
search (NSGA-III)~\cite{wang2023,wang2025nasv2}.

For Bayesian optimisation we used the TPE sampler through the Optuna
framework~\cite{akiba2019,bergstra2011}.
We ran 150 trials with the mean MAE on a 2D rectangle validation window as
the objective.
For NSGA-II~\cite{deb2002,blank2020pymoo} we used a population of 20 over
15 generations, optimising jointly for prediction accuracy and model size.
For NSGA-III~\cite{deb2014,blank2020pymoo} we used a population of 24 over
15 generations with reference points to maintain search diversity.

We found that the three strategies produced markedly different architecture
profiles.
Table~\ref{tab:arch} shows what each search returned.

\begin{table}[H]
\centering
\caption{PINN architecture profiles selected by each NAS strategy.}
\label{tab:arch}
\begin{tabular}{lcccc}
\toprule
Label & Strategy & Layers & Neurons/layer & Activation \\
\midrule
Bayesian/TPE & Bayesian + TPE & 5 & 151 & ReLU \\
NSGA-II      & NSGA-II        & 3 & 153 & tanh \\
NSGA-III     & NSGA-III       & 3 &  75 & tanh \\
\bottomrule
\end{tabular}
\end{table}

We observed that Bayesian/TPE returned the deepest and widest network at
approximately 93K trainable parameters in 3D mode.
We found that this architecture handled longer prediction windows better than
the other two throughout all experiments.
NSGA-II returned a compact three-layer design with roughly 48K parameters;
NSGA-III selected an even smaller network at around 12K parameters.
We observed that these lighter networks trained faster per window but required
shorter windows to remain stable, particularly on complex geometries where the
temperature gradient changes rapidly.

\begin{figure}[htbp]
  \centering
  \includegraphics[width=0.85\textwidth]{fig1_naspinn_2d.png}
  \caption{2D benchmark geometries from Wang et al.~\cite{wang2023} used as
  NAS validation and transfer targets.}
  \label{fig:naspinn_source_domains}
\end{figure}

We chose the 2D rectangle as the NAS validation domain because it is the
simplest benchmark: uniform boundary normals, no reentrant corners, and a
smooth temperature field.
An architecture that performs well here may not transfer perfectly to
geometries with curved surfaces or reentrant edges; evaluating that transfer
across all eight domains is one of the main results of the study
(Section~\ref{sec:disc_nas_transfer}).

\subsection{Training Protocol (Adam + L-BFGS, Per-Window Budget)}

Each PINN window was trained using two optimisers in sequence.
First, Adam \cite{kingma2015} was run with a cosine annealing schedule
\cite{loshchilov2017}, starting at $lr_0 = 8\times10^{-4}$ and decaying to
$lr_\mathrm{min} = 10^{-5}$ over $N_\mathrm{ep}$ epochs.
Once Adam converged, L-BFGS \cite{liu1989} was applied for a short
fine-tuning phase to sharpen the solution near the optimum.
This two-phase schedule gave lower final residuals than either optimiser
alone, particularly near the fin tip and reentrant corner regions where the
gradient is steepest.

The training budget varied by domain complexity.
Table~\ref{tab:training_budget} lists the configuration for each domain
group, together with the resulting per-window training time measured on a
single NVIDIA Tesla V100 GPU (full hardware and software specification in
Appendix~\ref{app:hardware}).

\begin{table}[H]
\centering
\caption{Training configuration and per-window wall-clock time by domain group
(NVIDIA Tesla V100). $^*$Bayesian/TPE on Thermal Fin uses Fourier embedding.}
\label{tab:training_budget}
\small
\begin{tabular}{llccccc}
\toprule
Domain group & Architecture & Adam epochs & L-BFGS & $n_\mathrm{domain}$ & $n_\mathrm{bc}$ &
  Time / window \\
\midrule
2D canonical  & Bayesian/TPE   &   800 &  50 & 1\,500 &   300 & $\sim 10\,\mathrm{s}$ \\
              & NSGA-II/III    & 1\,500 &  50 & 1\,500 &   300 & $\sim 16\,\mathrm{s}$ \\
\midrule
3D canonical  & Bayesian/TPE   &   800 & 100 & 12\,000 & 2\,400 & $\sim 25\,\mathrm{s}$ \\
              & NSGA-II/III    & 1\,500 & 100 & 12\,000 & 2\,400 & $\sim 35\,\mathrm{s}$ \\
\midrule
Thermal Fin   & Bayesian/TPE$^*$ & 2\,000 & 100 & 12\,000 & 2\,400 & $\sim 43\,\mathrm{s}$ \\
              & NSGA-II/III      & 2\,000 & 100 & 12\,000 & 2\,400 & $\sim 31\,\mathrm{s}$ \\
\bottomrule
\end{tabular}
\end{table}

The epoch counts differ between Bayesian/TPE (800) and the two evolutionary
architectures (1\,500) for the canonical domains.
The compact tanh networks of NSGA-II and NSGA-III require more gradient steps
to reach the same accuracy level that the larger ReLU network achieves in
fewer epochs.
Despite this difference, per-window training times remain in the same order of
magnitude across all three architectures for a given domain group.
We found that extending from 800 to 1\,500 epochs for NSGA-II/III improves
accuracy slightly but does not change the optimal $k$ on most domains.

One consequence of the window-by-window structure is that the total training
cost scales linearly with the number of windows, not with the window length.
Each window is trained independently from the FEM anchor field, so a
7.5-second window ($k=5$) takes the same time to train as a 1.5-second one
($k=1$).
At $k=5$ the Thermal Fin trains in 4 windows (around 2 minutes for
Bayesian/TPE); at $k=1$ the same 20 windows require around 10 minutes.
Once training is complete, inference at each window takes only milliseconds,
so the trained model can be reused across many simulation runs at negligible
marginal cost.

\section{Benchmark Domains}
\label{sec:domains}

We evaluated the framework on ten unique benchmark geometries.
The canonical study covers eight domains organised in three levels of increasing
complexity.
The adaptive-$k$ NAS architecture transfer study uses four 2D domains: Square,
Circle, L-Shape, and Flower.
Circle and L-Shape are shared with the canonical set; Square and Flower are
new domains that appear only in the adaptive-$k$ study.
Figure~\ref{fig:domains_overview} and Table~\ref{tab:domain_overview} present
the geometry and key properties of each domain together.

\begin{figure}[H]
\centering
\includegraphics[width=0.82\textwidth]{fig2_thesis_2d.png}
\subcaption{2D canonical domains}
\label{fig:domains_2d}

\vspace{0.8em}
\includegraphics[width=0.82\textwidth]{fig3_thesis_3d.png}
\subcaption{3D canonical domains}
\label{fig:domains_3d}

\vspace{0.8em}
\includegraphics[width=0.82\textwidth]{fig4_thermalfin.png}
\subcaption{Thermal Fin benchmark}
\label{fig:thermal_fin_domain}

\caption{Benchmark domain geometries.}
\label{fig:domains_overview}
\end{figure}

\begin{table}[H]
\centering\small\setlength{\tabcolsep}{4pt}
\caption{Benchmark domain summary: geometry, role, and key geometric challenge.
All domains use $T_0=540\,^\circ\mathrm{C}$, $T_w=20\,^\circ\mathrm{C}$, 30-second quench.}
\label{tab:domain_overview}
\begin{tabular}{@{}llp{3.5cm}p{6.5cm}@{}}
\toprule
\textbf{Domain} & \textbf{Dim.} & \textbf{Role} & \textbf{Geometric challenge / key property} \\
\midrule
\multicolumn{4}{l}{\textit{Level 1 — Adaptive-$k$ NAS Transfer \quad $h=5000\,\mathrm{W/(m^2K)}$}} \\
\midrule
Square 2D      & 2D & Regular geometry        & Uniform BC; no gradient concentration; simplest adaptive-$k$ test case \\
Circle 2D      & 2D & Curved boundary         & Curved normals; radial gradient; tests cut-cell BC \\
L-Shape 2D     & 2D & Reentrant corner        & 270° corner concentrates gradient; analogous to subframe thin rails \\
Flower 2D      & 2D & Concave boundary        & Multiple concave boundary regions; most geometrically complex 2D domain \\
\midrule
\multicolumn{4}{l}{\textit{Level 2 — Canonical 2D fixed-$k$ sweep \quad $h=5000\,\mathrm{W/(m^2K)}$ \quad (Circle and L-Shape shared with Level~1)}} \\
\midrule
Rectangle 2D   & 2D & NAS validation baseline & Simplest case; uniform BC; no gradient concentration \\
Circle 2D      & 2D & Curved boundary         & Shared with Level~1; fixed-$k$ sweep $k=1$--$5$ \\
L-Shape 2D     & 2D & Reentrant corner        & Shared with Level~1; fixed-$k$ sweep $k=1$--$5$ \\
\midrule
\multicolumn{4}{l}{\textit{Level 3 — Canonical 3D \quad $h=4000\,\mathrm{W/(m^2K)}$}} \\
\midrule
Rectangular 3D & 3D & 3D baseline             & Verifies 2D$\to$3D scaling; uniform cooling \\
Cylinder 3D    & 3D & Curved 3D surface       & Curved surface normals; azimuthal symmetry \\
Stacked 3D     & 3D & Material interface      & Two-layer gradient discontinuity; analogue of mould--part interface \\
L-Shape 3D     & 3D & 3D reentrant edge       & Edge gradient concentration in 3D; only domain where NSGA-III leads \\
\midrule
\multicolumn{4}{l}{\textit{Level 4 — Community benchmark \quad $h=4000\,\mathrm{W/(m^2K)}$, threshold $15\,^\circ\mathrm{C}$}} \\
\midrule
Thermal Fin    & 2D$\to$3D & External reference  & Original RBniCS benchmark is 2D steady-state~\cite{maday2004thermalfin}; adapted to 3D transient quenching in this thesis; most geometrically demanding case \\
\bottomrule
\end{tabular}
\end{table}


Full grid dimensions and active node counts for every domain are listed in
Appendix~\ref{app:fem_config}.

\subsection{Engineering Acceptance Threshold (MAE in $^\circ$C)}
\label{subsec:threshold}

The evaluation metric throughout this thesis is the mean-window MAE
($^\circ\mathrm{C}$), averaged over all prediction windows in the 30-second
simulation.
The objective is not to find the skip factor that minimises MAE, but to find
the maximum $k$ for which the MAE remains within an engineering acceptance
threshold.
This maximises FEM call reduction without compromising practical accuracy.

A percentage-of-range metric is used rather than an absolute error.
The physical consequence of a temperature prediction error depends on how far
that error is from the equilibrium state.
An error of $15\,^\circ\mathrm{C}$ near the quench temperature
($20\,^\circ\mathrm{C}$) matters far more than the same error when the part is
still near $540\,^\circ\mathrm{C}$.
The total quench range $\Delta T = T_0 - T_w = 520\,^\circ\mathrm{C}$ is
therefore used as the normalisation base.
Two domain-specific thresholds are defined.
For the Thermal Fin benchmark, the threshold is $15\,^\circ\mathrm{C}$
($2.9\%$ of $\Delta T$).
This value is derived from the accuracy requirements reported by
Mortensen et al.~\cite{mortensen2026}.
Residual stress and final dimension errors in that study remain within
forming tolerances when the transient temperature field is predicted to
within approximately $3\%$ of the operating range.
For the canonical 2D and 3D benchmarks no hard threshold is imposed.
Results are reported across the full $k=1$--$5$ range and the maximum
tolerable $k$ is discussed in the context of each domain.

The maximum acceptable $k$ per architecture is reported in each results
section alongside the corresponding FEM call reduction.
At $k=1$ the PINN covers one $1.5\,\mathrm{s}$ FEM step, which is the
baseline with no saving.
At $k=5$ it covers $7.5\,\mathrm{s}$, replacing four out of five FEM calls.
The practical contribution of this thesis is the mapping from geometry and
architecture to the maximum $k$ that keeps errors within the tolerance above.

\subsection{Physical Parameters}

All simulations use A356 aluminium material properties.
The same physical setup was used across all domain groups.
A slightly different value of $h$ was used for the 2D group, which was treated
as a separate preliminary benchmark.
Table~\ref{tab:params} lists the parameters used in each group.

\begin{table}[H]
\centering
\caption{Physical and simulation parameters.}
\label{tab:params}
\small
\begin{tabular}{lcc}
\toprule
Parameter & 2D group & 3D / Thermal Fin \\
\midrule
Thermal conductivity $k_T$     & $150\,\mathrm{W/(mK)}$                & $160\,\mathrm{W/(mK)}$ \\
Vol.\ heat capacity $\rho c_p$ & $2.4{\times}10^6\,\mathrm{J/(m^3K)}$  & $2.4{\times}10^6\,\mathrm{J/(m^3K)}$ \\
Thermal diffusivity $\alpha$   & $6.25{\times}10^{-5}\,\mathrm{m^2/s}$ & $6.67{\times}10^{-5}\,\mathrm{m^2/s}$ \\
Initial temperature $T_0$      & $540\,^\circ\mathrm{C}$                & $540\,^\circ\mathrm{C}$ \\
Water temperature $T_w$        & $20\,^\circ\mathrm{C}$                 & $20\,^\circ\mathrm{C}$ \\
Convection coefficient $h$     & $5000\,\mathrm{W/(m^2K)}$              & $4000\,\mathrm{W/(m^2K)}$ \\
Cooling time                   & $30\,\mathrm{s}$                       & $30\,\mathrm{s}$ \\
FEM internal step              & $0.5\,\mathrm{s}$                      & $0.5\,\mathrm{s}$ \\
FEM anchor spacing             & $1.5\,\mathrm{s}$                      & $1.5\,\mathrm{s}$ \\
\bottomrule
\end{tabular}
\end{table}

All parameters in Table~\ref{tab:params} are treated as temperature-independent
constants.
This simplifies the heat equation to a linear PDE and makes the PINN training
loss convex-like, at the cost of omitting the five-regime boiling curve and
temperature-dependent material properties used by Mortensen
et al.~\cite{mortensen2026} (described in
Section~\ref{sec:mortensen_reference}).
The rationale and impact of these simplifications are discussed in
Section~\ref{sec:conc_limitations}.
The implementation path for incorporating $h(T)$ is outlined in
Section~\ref{sec:conc_future}.


\section{Evaluation Metrics}
\label{sec:metrics}

Performance was measured using two error metrics and one efficiency metric.
For accuracy, the MAE in degrees Celsius is:

\begin{equation}
\mathrm{MAE} = \frac{1}{N}\sum_{i=1}^N |T_i^\mathrm{PINN} - T_i^\mathrm{FEM}|.
\label{eq:mae}
\end{equation}

This is computed over all prediction windows (mean-window MAE) and separately
for the last window only (final window MAE).
The mean-window MAE is the primary metric.
The final window value is misleading in PINN-only mode: as $t\to30\,\mathrm{s}$
the temperature field naturally converges toward $T_w=20\,^\circ\mathrm{C}$,
which reduces the absolute error regardless of how well the PINN actually
predicted the cooling history.

For full-field agreement, the relative $L_2$ error is:

\begin{equation}
\mathrm{Rel.}\;L_2
= \frac{\|\hat{T}-T^\mathrm{FEM}\|_2}{\|T^\mathrm{FEM}-T_w\|_2},
\label{eq:l2}
\end{equation}

where $T_w = 20\,^\circ\mathrm{C}$ is the quench temperature.
Normalising by $\|T^\mathrm{FEM}-T_w\|_2$ rather than $\|T^\mathrm{FEM}\|_2$
removes the constant baseline and measures error relative to the temperature
excursion above the quench medium, which is the physically meaningful driving
quantity.

For efficiency, the FEM anchor reduction $R(k)$ from
Eq.~\eqref{eq:reduction} is reported.
Since each anchor call corresponds to one FEM solver invocation, this metric
directly quantifies how much of the original FEM cost is saved.
The full implementation is provided in Appendix~\ref{app:code}.
